% StockSync Final Report — CS4391 Spring 2026
% Paste into Overleaf using the CVPR template (keep all \usepackage lines from template)
% Upload these figures to Overleaf:
%   results.png, confusion_matrix_normalized.png,
%   MaskPR_curve.png, val_batch0_pred.jpg, labels.jpg, train_batch0.jpg

\documentclass[10pt,twocolumn,letterpaper]{article}
\usepackage{cvpr}
\usepackage{times}
\usepackage{epsfig}
\usepackage{graphicx}
\usepackage{amsmath}
\usepackage{amssymb}
\usepackage{booktabs}
\usepackage[breaklinks=true,bookmarks=false]{hyperref}

\cvprfinalcopy

\def\cvprPaperID{****}
\def\httilde{\mbox{\tt\raisebox{-.5ex}{\symbol{126}}}}

\begin{document}

%%%%%%%%%%%%%%%%%%%%%%%%%%%%%%%%%%%%%%%%%%%%%%%%%%%%%%%%%%%%%%%%%%%%%%%%%%%%%%%%
\title{StockSync: Real-Time Automated Retail Inventory Tracking on\\
NVIDIA Jetson Orin Nano Using YOLOv8 Instance Segmentation}

\author{
Disha Shetty \quad Aashlesha [LASTNAME] \quad Hassan [LASTNAME] \quad Kashish [LASTNAME]\\
University of Texas at Dallas\\
{\tt\small \{[UTD NetIDs]\}@utdallas.edu}
}

\maketitle

%%%%%%%%%%%%%%%%%%%%%%%%%%%%%%%%%%%%%%%%%%%%%%%%%%%%%%%%%%%%%%%%%%%%%%%%%%%%%%%%
\begin{abstract}
We present StockSync, a real-time automated retail inventory
tracking system deployed on the NVIDIA Jetson Orin Nano.
The system uses a custom-trained YOLOv8n instance segmentation
model to detect and count shelf items from a live camera feed.
We collected and labeled a custom dataset of 425 images across
five product classes---chocolate, granola, mouse, pen, and
waterbottle---using Roboflow. The model was trained on Google
Colab with GPU acceleration and optimized using NVIDIA TensorRT
for edge deployment. Our model achieves a mask mAP@50 of
88.7\% on the held-out test set, with near-perfect performance
on waterbottle (98.5\%) and granola (99.2\%). The system
outputs real-time per-class inventory counts to a live
dashboard, enabling automated shelf monitoring without manual
inspection.
\end{abstract}

%%%%%%%%%%%%%%%%%%%%%%%%%%%%%%%%%%%%%%%%%%%%%%%%%%%%%%%%%%%%%%%%%%%%%%%%%%%%%%%%
\section{Introduction}

Retail inventory management is a labor-intensive process.
Employees must periodically walk store aisles to manually
count shelf stock---a process that is time-consuming,
error-prone, and does not scale to large stores. Out-of-stock
events cost retailers an estimated \$1 trillion annually in
lost sales globally~\cite{gruen2002retail}. Automating
inventory tracking through computer vision offers a practical
solution: a camera mounted above a shelf can continuously
monitor stock levels and alert staff when items run low.

We present StockSync, a shelf inventory tracking system that
runs on the NVIDIA Jetson Orin Nano developer kit---an
affordable, low-power edge AI platform suitable for
in-store deployment. Our system detects and segments
individual retail items in real time using a fine-tuned
YOLOv8n-seg model, counts per-class instances per frame, and
displays live inventory counts on a web dashboard.

The key contributions of this project are:
\begin{enumerate}
  \item A custom dataset of 425 shelf images across 5 product
        classes, collected with real cameras and labeled with
        instance segmentation masks.
  \item A fine-tuned YOLOv8n-seg model achieving 88.7\%
        mask mAP@50 on unseen test data.
  \item A TensorRT-optimized inference pipeline deployed on
        the NVIDIA Jetson Orin Nano for real-time edge
        inference.
  \item A full-stack dashboard displaying live per-class
        inventory counts from the camera feed.
\end{enumerate}

%%%%%%%%%%%%%%%%%%%%%%%%%%%%%%%%%%%%%%%%%%%%%%%%%%%%%%%%%%%%%%%%%%%%%%%%%%%%%%%%
\section{Related Work}

\noindent\textbf{Object Detection on Edge Devices.}
YOLO (You Only Look Once)~\cite{redmon2016yolo} introduced
single-stage real-time object detection. YOLOv8~\cite{yolov8}
extends this with an anchor-free architecture and native
support for instance segmentation, making it well-suited for
edge deployment. Several works have deployed YOLO variants on
NVIDIA Jetson platforms for applications such as traffic
monitoring and industrial inspection.

\noindent\textbf{Retail Inventory Tracking.}
Vision-based shelf monitoring has been explored using both
classical computer vision and deep learning approaches.
Recent systems use object detection to count items and detect
gaps on shelves~\cite{goldman2019precise}. Our work differs
in combining instance segmentation with full edge deployment
on low-cost embedded hardware using a fully custom dataset.

\noindent\textbf{Instance Segmentation.}
Mask R-CNN~\cite{he2017mask} introduced the standard
two-stage approach to instance segmentation. YOLOv8-seg
adapts this to a single-stage real-time framework using a
lightweight mask head, enabling deployment on constrained
hardware such as the Jetson Orin Nano without sacrificing
segmentation quality.

\noindent\textbf{TensorRT Optimization.}
NVIDIA TensorRT~\cite{tensorrt} optimizes deep learning
models for inference on NVIDIA GPUs through layer fusion,
precision calibration (FP16/INT8), and kernel auto-tuning.
Prior work demonstrates 2--5$\times$ speedup on Jetson
platforms compared to standard PyTorch inference, enabling
real-time performance on edge hardware.

%%%%%%%%%%%%%%%%%%%%%%%%%%%%%%%%%%%%%%%%%%%%%%%%%%%%%%%%%%%%%%%%%%%%%%%%%%%%%%%%
\section{Method}

\subsection{System Overview}

The StockSync pipeline consists of four stages: (1) data
collection and labeling, (2) model training on cloud GPU,
(3) TensorRT optimization, and (4) Jetson deployment with
live dashboard. A camera feeds frames into the inference
pipeline, which outputs per-class instance counts to a
backend server and live web dashboard.

\subsection{Dataset}

We collected 425 images using a Logitech C270 HD Webcam and
an 8MP IMX219 camera across five product classes: chocolate,
granola, mouse, pen, and waterbottle. Images were captured
under varying lighting conditions and shelf arrangements to
improve generalization. Annotation was performed using
Roboflow with instance segmentation (polygon) masks.

\begin{table}[h]
\centering
\small
\begin{tabular}{lrrrr}
\toprule
\textbf{Class} & \textbf{Train} & \textbf{Val} &
\textbf{Test} & \textbf{Total} \\
\midrule
waterbottle & 113 & 29 & 23 & 165 \\
granola     &  87 & 15 & 12 & 114 \\
chocolate   &  68 & 21 & 16 & 105 \\
pen         &  33 &  8 &  9 &  50 \\
mouse       &  18 &  2 &  2 &  22 \\
\midrule
\textbf{Total} & \textbf{319} & \textbf{75} &
\textbf{62} & \textbf{456} \\
\bottomrule
\end{tabular}
\caption{Images per class per split. Counts reflect unique
images containing at least one instance of each class.}
\label{tab:dataset}
\end{table}

Figure~\ref{fig:dataset} shows the class instance
distribution and bounding box statistics. Waterbottle is the
most represented class with 341 training instances, while
mouse is the most underrepresented with only 30.

\begin{figure}[h]
  \centering
  \includegraphics[width=\linewidth]{labels.jpg}
  \caption{Dataset statistics: class instance distribution
  (top left), bounding box spatial distribution (bottom left
  and right), and aspect ratio distribution (top right).}
  \label{fig:dataset}
\end{figure}

Figure~\ref{fig:trainbatch} shows a sample training batch
with ground-truth segmentation masks overlaid.

\begin{figure}[h]
  \centering
  \includegraphics[width=\linewidth]{train_batch0.jpg}
  \caption{Sample training batch showing instance
  segmentation masks across multiple product classes.}
  \label{fig:trainbatch}
\end{figure}

\subsection{Model: YOLOv8n-seg}

We selected YOLOv8n-seg (nano segmentation variant) for its
balance of accuracy and inference speed on constrained
hardware. The model uses a CSP (Cross Stage Partial)
backbone, PANet neck, and a lightweight segmentation head
that predicts both bounding boxes and polygon masks. With
3.26M parameters and 11.5 GFLOPs, it is well-suited for
real-time inference on the Jetson Orin Nano.

We initialized from ImageNet-pretrained weights
(\texttt{yolov8n-seg.pt}) and fine-tuned on our custom
dataset, overriding the output head for $N_c = 5$ classes.
381 of 417 pretrained weight tensors were transferred
successfully.

\subsection{Training}

Training was performed on Google Colab with an NVIDIA Tesla
T4 GPU (16\,GB VRAM). Key hyperparameters are listed in
Table~\ref{tab:hyperparams}.

\begin{table}[h]
\centering
\small
\begin{tabular}{ll}
\toprule
\textbf{Hyperparameter} & \textbf{Value} \\
\midrule
Max epochs     & 100 (early stopped at 65) \\
Best epoch     & 50 \\
Batch size     & 8 \\
Image size     & 640$\times$640 \\
Optimizer      & AdamW \\
Learning rate  & 0.001111 \\
Momentum       & 0.9 \\
Weight decay   & 0.0005 \\
Early stopping & patience = 15 \\
Training time  & $\approx$10 minutes \\
\bottomrule
\end{tabular}
\caption{Training hyperparameters.}
\label{tab:hyperparams}
\end{table}

Data augmentation included random horizontal flip, HSV
color jitter, mosaic augmentation, random erasing, Gaussian
blur, and CLAHE contrast enhancement. The optimizer and
learning rate were determined automatically by Ultralytics'
\texttt{auto} optimizer selection.

Figure~\ref{fig:training} shows training and validation
loss curves alongside mAP metrics across epochs. All losses
decrease steadily and mAP plateaus around epoch 50, after
which early stopping triggered at epoch 65.

\begin{figure}[h]
  \centering
  \includegraphics[width=\linewidth]{results.png}
  \caption{Training curves: box loss, segmentation loss,
  and classification loss (top row: train, bottom row: val),
  along with precision, recall, mAP@50, and mAP@50-95
  metrics across 65 epochs.}
  \label{fig:training}
\end{figure}

\subsection{TensorRT Optimization}

The trained \texttt{best.pt} PyTorch model was exported to
TensorRT engine format on the Jetson Orin Nano using the
Ultralytics export pipeline with FP16 precision. TensorRT
performs layer fusion, memory optimization, and kernel
auto-tuning for the Jetson's CUDA cores, resulting in
significantly faster inference compared to standard PyTorch.

\subsection{Deployment}

The TensorRT-optimized model runs on the Jetson Orin Nano
connected to the IMX219 camera. The inference pipeline
captures frames, runs detection, and counts per-class
instances. Results are sent to a backend server which feeds
a live web dashboard displaying real-time inventory counts
per product class, enabling store staff to monitor shelf
stock without manual inspection.

%%%%%%%%%%%%%%%%%%%%%%%%%%%%%%%%%%%%%%%%%%%%%%%%%%%%%%%%%%%%%%%%%%%%%%%%%%%%%%%%
\section{Experiments}

\subsection{Evaluation Metrics}

We evaluate using standard COCO metrics: mAP@50 (mean
Average Precision at IoU threshold 0.50) and mAP@50-95
(mean AP averaged over IoU thresholds 0.50--0.95). We
report both box (detection) and mask (segmentation) metrics.
Final evaluation uses the held-out test set (60 images,
178 instances) not seen during training or validation.

\subsection{Quantitative Results}

Table~\ref{tab:results} reports overall validation and test
set results. The model generalizes well, with only a modest
drop from val (mask mAP@50 = 0.944) to test (0.887),
indicating limited overfitting despite the small dataset.

\begin{table}[h]
\centering
\small
\begin{tabular}{lcccc}
\toprule
& \multicolumn{2}{c}{\textbf{Box}} &
  \multicolumn{2}{c}{\textbf{Mask}} \\
\cmidrule(lr){2-3}\cmidrule(lr){4-5}
\textbf{Split} & mAP@50 & mAP@50-95 & mAP@50 & mAP@50-95 \\
\midrule
Val  & 0.945 & 0.843 & 0.944 & 0.814 \\
Test & 0.913 & 0.784 & 0.887 & 0.772 \\
\bottomrule
\end{tabular}
\caption{Overall validation and test set results.}
\label{tab:results}
\end{table}

Table~\ref{tab:perclass} shows per-class test results.
Waterbottle and granola achieve near-perfect segmentation,
consistent with their larger training sets. Pen is the
weakest class at 0.601 mAP@50, attributed to its small,
thin geometry and limited training data (33 images).

\begin{table}[h]
\centering
\small
\begin{tabular}{lcccc}
\toprule
\textbf{Class} & \textbf{P} & \textbf{R} &
\textbf{mAP@50} & \textbf{mAP@50-95} \\
\midrule
waterbottle & 1.000 & 0.983 & 0.985 & 0.948 \\
mouse       & 0.618 & 1.000 & 0.995 & 0.920 \\
granola     & 0.975 & 0.944 & 0.992 & 0.899 \\
chocolate   & 0.883 & 0.756 & 0.862 & 0.665 \\
pen         & 0.839 & 0.583 & 0.601 & 0.425 \\
\midrule
\textbf{All} & 0.867 & 0.858 & 0.887 & 0.772 \\
\bottomrule
\end{tabular}
\caption{Per-class mask results on held-out test set.}
\label{tab:perclass}
\end{table}

\subsection{Confusion Matrix Analysis}

Figure~\ref{fig:confusion} shows the normalized confusion
matrix on the validation set. Mouse, pen, granola, and
waterbottle are classified with high accuracy (1.00, 0.96,
0.94, 0.97 respectively). Chocolate has the most confusion,
with 37\% of instances missed as background---likely due to
visual similarity between chocolate packaging and background
surfaces, combined with the smaller number of training
images (68 images vs.\ 113 for waterbottle).

\begin{figure}[h]
  \centering
  \includegraphics[width=\linewidth]{confusion_matrix_normalized.png}
  \caption{Normalized confusion matrix on the validation
  set. Diagonal values indicate correctly classified
  fractions. Chocolate shows the most background confusion
  (0.37).}
  \label{fig:confusion}
\end{figure}

\subsection{Precision-Recall Analysis}

Figure~\ref{fig:pr} shows the mask Precision-Recall curve
per class. Mouse, waterbottle, and granola maintain high
precision across the full recall range. Pen drops off
earliest, reflecting its difficulty. The overall mask
mAP@50 = 0.943 on the validation set confirms strong
aggregate performance.

\begin{figure}[h]
  \centering
  \includegraphics[width=\linewidth]{MaskPR_curve.png}
  \caption{Mask Precision-Recall curves per class on the
  validation set. Higher area under the curve indicates
  better performance.}
  \label{fig:pr}
\end{figure}

\subsection{Qualitative Results}

Figure~\ref{fig:preds} shows sample inference outputs from
the validation set. The model produces accurate segmentation
masks with high confidence scores across multiple product
classes in varied backgrounds.

\begin{figure}[h]
  \centering
  \includegraphics[width=\linewidth]{val_batch0_pred.jpg}
  \caption{Sample model predictions on validation images,
  showing segmentation masks and confidence scores.}
  \label{fig:preds}
\end{figure}

\subsection{TensorRT Inference Speed}

Table~\ref{tab:speed} reports inference speed on the Jetson
Orin Nano comparing standard PyTorch and TensorRT FP16.

\begin{table}[h]
\centering
\small
\begin{tabular}{lcc}
\toprule
\textbf{Mode} & \textbf{ms/frame} & \textbf{FPS} \\
\midrule
PyTorch (CPU)    & [X] & [X] \\
PyTorch (GPU)    & [X] & [X] \\
TensorRT FP16    & [X] & [X] \\
\bottomrule
\end{tabular}
\caption{Inference speed on NVIDIA Jetson Orin Nano.
TensorRT FP16 provides significant speedup over PyTorch,
enabling real-time performance.}
\label{tab:speed}
\end{table}

\subsection{Analysis}

Performance correlates strongly with training set size
(Table~\ref{tab:dataset}). The three best-performing
classes (waterbottle, granola, mouse) benefit from
distinctive visual features, while pen suffers from both
limited data and inherent geometric difficulty---thin
objects with high aspect ratio produce small segmentation
masks that are harder to localize. Mouse achieves high
mAP@50 (0.995) but was evaluated on only 2 test images,
so this result should be interpreted cautiously.

The modest val-to-test drop (0.944$\rightarrow$0.887)
suggests the model generalizes reasonably well despite the
small dataset size. The primary improvement avenue is
additional data collection for pen ($\sim$50 more images)
and mouse ($\sim$80 more images).

%%%%%%%%%%%%%%%%%%%%%%%%%%%%%%%%%%%%%%%%%%%%%%%%%%%%%%%%%%%%%%%%%%%%%%%%%%%%%%%%
\section{Conclusion}

We presented StockSync, a real-time retail shelf inventory
tracking system running on the NVIDIA Jetson Orin Nano. By
fine-tuning YOLOv8n-seg on a custom 425-image dataset, we
achieved 88.7\% mask mAP@50 on unseen test data. Three of
five classes exceed 98\% mAP@50, demonstrating the
viability of edge-deployed instance segmentation for retail
monitoring. Performance variation across classes correlates
directly with training data size, confirming that data
collection is the primary bottleneck for improvement.

The full pipeline---data collection, annotation, cloud
training, TensorRT optimization, and edge deployment with a
live dashboard---demonstrates a complete end-to-end computer
vision system suitable for real-world retail inventory
applications. Future work should prioritize collecting
additional images for pen and mouse, and explore model
distillation or quantization to further reduce latency on
the Jetson.

%%%%%%%%%%%%%%%%%%%%%%%%%%%%%%%%%%%%%%%%%%%%%%%%%%%%%%%%%%%%%%%%%%%%%%%%%%%%%%%%
\section*{Contributions}

\noindent\textbf{Disha Shetty} served as team lead and
was responsible for overall project coordination, task
assignment, and technical direction. Collected 153
waterbottle images for the dataset. Authored the project
proposal, mid-term report, and final report. Set up and
validated the NVIDIA Jetson Orin Nano hardware environment
and recorded all required video submissions demonstrating
model execution on the device. Designed and implemented
the complete ML pipeline: selected YOLOv8n-seg as the
model architecture, configured and executed training on
Google Colab with GPU acceleration, performed
hyperparameter tuning, and conducted validation and
test set evaluation achieving 88.7\% mask mAP@50.
Implemented TensorRT optimization for edge deployment
and developed the real-time inference pipeline with live
segmentation mask overlay and per-class inventory count
display for the Jetson demo.

\noindent\textbf{Hassan [LASTNAME]} collected granola bar
images for the dataset. Designed and implemented the
backend data pipeline, including an SQLite-based detection
logging system, a CSV export layer for frame-level
inventory counts, and a FastAPI server exposing REST
endpoints for querying detection history and session
summaries. Validated the backend pipeline end-to-end by
running inference sessions and verifying data storage and
retrieval across multiple test runs.

\noindent\textbf{Aashlesha [LASTNAME]} collected 30 mouse
and 48 pen images for the dataset and led the data
annotation effort for all five product classes using
Roboflow, producing instance segmentation polygon masks
used to train the model.

\noindent\textbf{Kashish [LASTNAME]} collected chocolate
images for the dataset. Was responsible for developing
the frontend dashboard and ensuring it was fully
integrated with the backend system. Designed and
implemented the user interface to display real-time
inventory data in a clean and intuitive way, including
metrics, live detection feed, and inventory counts. Once
the backend APIs were completed, connected the frontend
to these endpoints, enabling seamless data flow and
real-time updates. This integration allowed the system
to function as a cohesive application, transforming
backend outputs into a user-friendly and interactive
dashboard.

{\small
\bibliographystyle{ieee_fullname}
\begin{thebibliography}{9}

\bibitem{redmon2016yolo}
J.~Redmon, S.~Divvala, R.~Girshick, and A.~Farhadi.
\newblock You only look once: Unified, real-time object detection.
\newblock In {\em CVPR}, 2016.

\bibitem{he2017mask}
K.~He, G.~Gkioxari, P.~Doll{\'a}r, and R.~Girshick.
\newblock Mask {R-CNN}.
\newblock In {\em ICCV}, 2017.

\bibitem{yolov8}
G.~Jocher, A.~Chaurasia, and J.~Qiu.
\newblock {Ultralytics YOLOv8}.
\newblock \url{https://github.com/ultralytics/ultralytics}, 2023.

\bibitem{tensorrt}
NVIDIA.
\newblock {TensorRT} developer guide.
\newblock \url{https://docs.nvidia.com/deeplearning/tensorrt/}, 2023.

\bibitem{gruen2002retail}
T.~W. Gruen, D.~S. Corsten, and S.~Bharadwaj.
\newblock Retail out-of-stocks: A worldwide examination of extent, causes, and
  consumer responses.
\newblock {\em Grocery Manufacturers of America}, 2002.

\bibitem{goldman2019precise}
S.~Goldman, Y.~Reich, D.~Tsur, and E.~Wiesel.
\newblock Precise detection in densely packed scenes.
\newblock In {\em CVPR}, 2019.

\end{thebibliography}
}

\end{document}
